\chapter*{Abstract}
The compiler is an essential tool in the development of computer programs that
make efficient use of the system's hardware resources. It is often treated,
however, as a mere code translator, with its optimization capabilities left
aside, even though the way source code is translated into machine instructions
directly affects program performance.

This issue is of particular importance in embedded systems, where improved
performance often means lower cost and reduced power consumption. This thesis
presents an embedded system based on the Raspberry Pi 4B platform running Linux,
designed to study the optimization capabilities of the GCC and Clang compilers.
The fast Fourier transform was chosen as the test algorithm, for which eleven
variants were implemented in C, both sequential and multithreaded.

The thesis also covers operating systems in embedded applications, the
fundamentals of the Fourier transform, and the structure of a compiler together
with the optimizations it applies. The study covered eleven sets of compilation
flags and three input data sizes; in addition to execution time, hardware
counters, output code size, and computational accuracy were analysed.

The largest gain came from moving away from the default \texttt{-O0} level,
which reduced execution time by a factor of 2.18 to 5.11, whereas further flag
selection typically yielded around 12\%. For large data sets, execution time was
determined mainly by the memory access pattern, which made the choice of
computational scheme as important as the optimization level. Manual
vectorization using NEON instructions and parallelization of the computation
brought clear benefits only within specific ranges of data size. The compiler
did not correct numerical errors resulting from the construction of the
algorithm. The results indicate that program optimization should begin with the
choice of algorithm and data organization, with compilation settings treated as
a complement to them.

\bigskip

\noindent\textbf{Keywords:} compilation optimization, GCC and Clang compilers,
fast Fourier transform, embedded systems, Linux operating system

\bigskip

\noindent\textbf{OECD consistent field of science and technology classification:} Engineering and technology, Electrical engineering, Electronic engineering, Information engineering, 
